$resume-preamble.tex()$

\begin{document}

%-------------------------------------------------------------------------------
%	HEADER
%-------------------------------------------------------------------------------
% LOCAL: resume.cls's \name takes one string rather than awesome-cv.cls's
% \name{firstname}{lastname}, so `surname` (when given) is folded into it here
% rather than exposed as a second class argument. Documented in the README.
\name{$name$$if(surname)$ $surname$$endif$}

$if(position)$
% No upstream equivalent -- upstream's own examples put a role only in
% \role{} beside an entry, never under the name. Set plainly, the same size as
% the contact line it sits above, so it reads as a subtitle rather than an entry.
{\centering\normalsize $position$\par}
\vspace{0.5ex}
$endif$

% Built from whichever of phone/email/www/github/linkedin the document sets,
% each drawn through resume.cls's fontawesome5-based macros and separated by
% \rescontactsep -- see the LOCAL note on \rescontact in resume.cls.
\basicInfo{%
$if(phone)$\rescontact{\phone{$phone$}}$endif$
$if(email)$\rescontact{\email{$email$}}$endif$
$if(www)$\rescontact{\homepage{$www$}}$endif$
$if(github)$\rescontact{\github{$github$}}$endif$
$if(linkedin)$\rescontact{\linkedin{$linkedin$}}$endif$
}

%-------------------------------------------------------------------------------
%	RESUME CONTENT
%	Each section is written as a heading in the .qmd; resume.cls's own
%	\titleformat restyles \section/\subsection directly (see README), so no
%	heading-remapping filter is needed here the way awesome-cv.lua provides one.
%-------------------------------------------------------------------------------
$if(linestretch)$
\setstretch{$linestretch$}
$endif$

$for(include-before)$
$include-before$
$endfor$

$body$

$before-bib.tex()$

$if(nocite-ids)$
\nocite{$for(nocite-ids)$$it$$sep$, $endfor$}
$endif$
$biblio.tex()$

$for(include-after)$
$include-after$
$endfor$

$after-body.tex()$

\end{document}
